\chapter{Complete Implementation of ML-DSA (FIPS 204)}
\chaptermark{ML-DSA (FIPS 204)}
\label{ch:mldsa}

\section{Learning objectives}

The previous part of the book built up to FIPS 203, the real ML-KEM standard for \emph{encryption and key exchange}. We now cross to the other half of public-key cryptography, \emph{digital signatures}, and to the first of three signature standards: FIPS 204, \textbf{ML-DSA}, better known by its submission name Dilithium~\cite{fips204,dilithium2018}.

After finishing this chapter, you should be able to:

\begin{enumerate}
  \item explain what a signature scheme does and how it differs from a KEM;
  \item follow the ML-DSA key generation flow and see that it reuses the module-lattice machinery of Kyber almost verbatim;
  \item explain the two hard problems, Module-LWE and Module-SIS, and which part of the scheme each one protects;
  \item follow the signing procedure phase by phase, and understand why it contains a rejection (``abort'') loop;
  \item connect the rejection idea back to the role of noise studied in the LWE chapters.
\end{enumerate}

\section{Review: how signing differs from ML-KEM}

ML-KEM answered the question \emph{how do two parties agree on a shared secret?} A signature answers a different question: \emph{how does a verifier know that a message really came from the holder of a secret key, unchanged?} The three algorithms mirror a KEM but serve this new goal:

\begin{itemize}
  \item \textbf{KeyGen} $\to$ a public verification key $pk$ and a secret signing key $sk$;
  \item \textbf{Sign}$(sk, M)$ $\to$ a signature $\sigma$ on a message $M$;
  \item \textbf{Verify}$(pk, M, \sigma)$ $\to$ \emph{accept} or \emph{reject}.
\end{itemize}

The security goal is \emph{unforgeability}: even after seeing many valid signatures, nobody without $sk$ can produce a valid signature on a new message. The remarkable fact is that ML-DSA achieves this on \emph{exactly} the same ring $R_q = \mathbb{Z}_q[x]/(x^{256}+1)$ and module structure we already built for Kyber. Only the modulus changes ($q = 8380417$), and the signing procedure is genuinely new.

\section{The two hard problems}

ML-DSA leans on two lattice problems in the module setting of Chapter~\ref{ch:mlwe}.

\begin{itemize}
  \item \textbf{Module-LWE}, already familiar: from $(A,\ \mathbf{t}=A\mathbf{s}_1+\mathbf{s}_2)$ with short $\mathbf{s}_1,\mathbf{s}_2$, recovering the secret is hard. This hides the signing key inside the public key.
  \item \textbf{Module-SIS} (Short Integer Solution): the module-lattice form of the SIS problem introduced in Section~\ref{sec:sis}. Given a random $A$, finding a short nonzero $\mathbf{z}$ with $A\mathbf{z}=\mathbf{0}\pmod q$ is hard. This blocks forgery: a forger would have to produce a short vector satisfying a public equation.
\end{itemize}

SIS is the mirror image of LWE. LWE hides a secret inside noise; SIS asks you to \emph{find} a short solution to a random system. Both are hard for the same geometric reason -- short vectors in a high-dimensional lattice are hard to find.

\section{The real input in FIPS 204}

Exactly as in FIPS 203, key generation does not begin by sampling a matrix. It begins with a single short seed:
\[
\xi \in \{0,1\}^{256}.
\]
Everything else -- the public matrix, both secret vectors, and the internal PRF keys -- is expanded deterministically from $\xi$. Key generation is therefore a deterministic function $\mathrm{KeyGen}(\xi)$, just like ML-KEM.

\section{Phase 1: expand the seed}

The seed is hashed by $H$ (built on SHAKE-256) and split into three parts:
\[
(\rho,\ \rho',\ K) = H(\xi).
\]
Their roles are:
\begin{itemize}
  \item $\rho$ -- public, expands into the matrix $A$ (it will be stored in the public key);
  \item $\rho'$ -- secret, seeds the sampling of $\mathbf{s}_1$ and $\mathbf{s}_2$;
  \item $K$ -- secret, a PRF key used later to derive per-signature randomness.
\end{itemize}

\section{Phase 2: generate the matrix}

As in Kyber, the matrix is never stored; it is regenerated on demand from its seed:
\[
A = \textsc{ExpandA}(\rho) \in R_q^{k\times \ell}.
\]
The dimensions $(k,\ell)$ are what distinguish the three security levels.

\section{Phase 3: sample the secret vectors}

The two short secrets are expanded from $\rho'$:
\[
(\mathbf{s}_1,\ \mathbf{s}_2) = \textsc{ExpandS}(\rho'),\qquad \mathbf{s}_1\in R_q^\ell,\ \mathbf{s}_2\in R_q^k,
\]
with coefficients in the small range $[-\eta,\eta]$. These play the role of the LWE secret and error.

\section{Phase 4: compute and split the public value}

The public value is the familiar Module-LWE equation:
\[
\mathbf{t} = A\mathbf{s}_1 + \mathbf{s}_2.
\]
To shrink the public key, $\mathbf{t}$ is split into a high part $\mathbf{t}_1$ (published) and a low part $\mathbf{t}_0$ (kept in the secret key):
\[
\mathbf{t} = \mathbf{t}_1\cdot 2^{d} + \mathbf{t}_0.
\]
The public key is $(\rho, \mathbf{t}_1)$; the secret key bundles $(\rho, K, \mathbf{t}_0, \mathbf{s}_1, \mathbf{s}_2)$ together with a hash $tr = H(pk)$ used during signing.

\begin{figure}[H]
\centering
\begin{tikzpicture}[every node/.style={font=\scriptsize}]
  \node[flownode] (xi) at (5,7.4) {seed $\xi \in \{0,1\}^{256}$};
  \node[flownode] (h)  at (5,6.2) {$H(\xi)$};
  \node[flownode] (rho)   at (1.3,5.0) {$\rho$};
  \node[flownode] (rhop)  at (5,5.0) {$\rho'$};
  \node[flownode] (K)     at (8.7,5.0) {$K$ (PRF key)};
  \node[flownode, font=\scriptsize] (A)     at (1.5,3.6) {$A=\textsc{ExpandA}(\rho)$};
  \node[flownode, font=\scriptsize] (s)     at (5.9,3.6) {$(\mathbf{s}_1,\mathbf{s}_2)=\textsc{ExpandS}(\rho')$};
  \node[flownode] (t)     at (3.6,2.4) {$\mathbf{t}=A\mathbf{s}_1+\mathbf{s}_2$};
  \node[flownode] (split) at (3.6,1.3) {split $\mathbf{t}=\mathbf{t}_1 2^d+\mathbf{t}_0$};
  \node[keynode]  (pk)    at (1.7,-0.1) {public key $(\rho,\mathbf{t}_1)$};
  \node[keynode]  (sk)    at (7.0,-0.1) {secret key\\[-2pt]$(\rho,K,tr,\mathbf{t}_0,\mathbf{s}_1,\mathbf{s}_2)$};
  \draw[flowarrow] (xi) -- (h);
  \draw[flowarrow] (h) -- (rho);
  \draw[flowarrow] (h) -- (rhop);
  \draw[flowarrow] (h) -- (K);
  \draw[flowarrow] (rho) -- (A);
  \draw[flowarrow] (rhop) -- (s);
  \draw[flowarrow] (A) -- (t);
  \draw[flowarrow] (s) -- (t);
  \draw[flowarrow] (t) -- (split);
  \draw[flowarrow] (split) -- (pk);
  \draw[flowarrow] (split) -- (sk);
  \draw[flowarrow] (K) to[out=270,in=0] (sk.east);
\end{tikzpicture}
\caption{ML-DSA key generation, in the same deterministic-expansion style as FIPS 203. A single seed $\xi$ fans out into the matrix, both secrets, and the PRF key; the public key stores only $(\rho,\mathbf{t}_1)$, from which $A$ can always be regenerated.}
\label{fig:mldsa-keygen}
\end{figure}

\begin{algorithm}[H]
\caption{ML-DSA.KeyGen (teaching form)}
\label{alg:mldsa-keygen}
\begin{algorithmic}[1]
\Require Seed $\xi \in \{0,1\}^{256}$; parameters $(k,\ell,\eta,d)$
\Ensure Public key $(\rho,\mathbf{t}_1)$; secret key $(\rho,K,tr,\mathbf{t}_0,\mathbf{s}_1,\mathbf{s}_2)$
\State $(\rho,\rho',K) \gets H(\xi)$ \Comment{Phase 1}
\State $A \gets \textsc{ExpandA}(\rho)$ \Comment{Phase 2}
\State $(\mathbf{s}_1,\mathbf{s}_2) \gets \textsc{ExpandS}(\rho')$ \Comment{Phase 3}
\State $\mathbf{t} \gets A\mathbf{s}_1 + \mathbf{s}_2$ \Comment{Phase 4}
\State $(\mathbf{t}_1,\mathbf{t}_0) \gets \textsc{Power2Round}(\mathbf{t}, d)$
\State $tr \gets H(\rho\,\|\,\mathbf{t}_1)$
\State \Return $\big((\rho,\mathbf{t}_1),\ (\rho,K,tr,\mathbf{t}_0,\mathbf{s}_1,\mathbf{s}_2)\big)$
\end{algorithmic}
\end{algorithm}

\section{Interlude: the Fiat--Shamir transform}
\label{sec:fiat-shamir}

Before diving into ML-DSA's signing, it is worth isolating the general technique it uses, because it is a basic building block that appears far beyond this one scheme.

An \textbf{identification scheme} lets a prover convince a verifier that it knows a secret $s$ (whose public image is $t$), without revealing $s$. It is a three-move conversation: the prover sends a \emph{commitment} $w$; the verifier replies with a random \emph{challenge} $c$; the prover sends a \emph{response} $z$ that only someone knowing $s$ could produce, and the verifier checks a relation among $(w,c,z,t)$. If the response is properly randomized, the transcript reveals nothing about $s$.

The \textbf{Fiat--Shamir transform}~\cite{fiatshamir1986} turns this interactive protocol into a signature by \emph{removing the verifier}. Instead of waiting for a random challenge, the prover computes the challenge itself as a hash of the commitment together with the message:
\[
c = H(w \,\|\, m).
\]
Because $H$ behaves like an unpredictable random function, the prover cannot steer $c$ to its advantage; and because $c$ depends on $m$, the whole transcript is bound to that message -- it \emph{is} a signature on $m$.

\begin{figure}[H]
\centering
\begin{tikzpicture}[font=\scriptsize]
  \node[font=\small\bfseries] at (1.5,3.5) {Interactive identification};
  \node[flownode] (P) at (0,2.9) {Prover};
  \node[flownode] (V) at (3,2.9) {Verifier};
  \draw[pqcgray!45, dashed] (0,2.5) -- (0,0.1);
  \draw[pqcgray!45, dashed] (3,2.5) -- (3,0.1);
  \draw[flowarrow] (0,2.0) -- node[above]{commitment $w$} (3,2.0);
  \draw[flowarrow] (3,1.35) -- node[above]{random $c$} (0,1.35);
  \draw[flowarrow] (0,0.7) -- node[above]{response $z$} (3,0.7);
  \node[pqcteal] at (3,0.25) {check};
  \node[pqcorange, font=\Large] at (4.7,1.6) {$\Rightarrow$};
  \node[font=\small\bfseries] at (8.2,3.5) {Non-interactive (Fiat--Shamir)};
  \node[flownode, align=center] (P2) at (8.2,2.3) {Prover computes $w$,\\then $c=H(w\,\|\,m)$ itself};
  \node[keynode] (sig) at (8.2,1.0) {signature $\sigma=(c,z)$};
  \draw[flowarrow] (P2) -- (sig);
\end{tikzpicture}
\caption{The Fiat--Shamir transform. On the left, an interactive identification protocol: three messages, with the verifier supplying a fresh random challenge $c$. On the right, the same transcript made non-interactive by having the prover compute $c=H(w\,\|\,m)$ itself -- the hash plays the role of the verifier, and binding it to the message $m$ turns the transcript into a signature. This recipe underlies many signature schemes; ML-DSA instantiates it with module lattices.}
\label{fig:fiat-shamir}
\end{figure}

\section{The signing idea: commitment, challenge, response}

Signing instantiates exactly that identification transcript, with the module-lattice objects playing the roles of commitment, challenge, and response:

\begin{enumerate}
  \item \textbf{Commit.} Sample a random short masking vector $\mathbf{y}$ and form $\mathbf{w} = A\mathbf{y}$.
  \item \textbf{Challenge.} First form the message representative $\mu=H(tr\,\|\,M')$, where $M'$ includes any context or pre-hash formatting. Then compute $\tilde c=H(\mu\,\|\,\mathbf{w}_1)$ with $\mathbf{w}_1=\textsc{HighBits}(\mathbf{w})$, and derive $c=\textsc{SampleInBall}(\tilde c)$.
  \item \textbf{Respond.} Reply with $\mathbf{z} = \mathbf{y} + c\,\mathbf{s}_1$.
\end{enumerate}

A verifier can then recompute $A\mathbf{z} - c\,\mathbf{t}$ and check it matches $\mathbf{w}$ in the high bits, because
\[
A\mathbf{z} - c\,\mathbf{t} = A(\mathbf{y}+c\,\mathbf{s}_1) - c(A\mathbf{s}_1+\mathbf{s}_2) = A\mathbf{y} - c\,\mathbf{s}_2 \approx \mathbf{w}.
\]
The leftover $c\,\mathbf{s}_2$ is small, and the \textsc{HighBits} rounding absorbs it. Replacing the interactive verifier by the hash $H$ is precisely the Fiat--Shamir transform: the hash acts as an unpredictable, unmanipulable challenger.

\section{Why signing must sometimes abort}

The response $\mathbf{z} = \mathbf{y} + c\,\mathbf{s}_1$ carries the secret $\mathbf{s}_1$ directly. If the mask $\mathbf{y}$ is not large enough to swamp $c\,\mathbf{s}_1$, then $\mathbf{z}$ leaks information -- across many signatures the mask averages out and the secret emerges. This is the LWE lesson in a new costume: \emph{the mask must completely hide the secret}.

ML-DSA's answer is \textbf{rejection sampling}, also called Fiat--Shamir \emph{with aborts}. After forming $\mathbf{z}$, the signer checks whether it lands in a safe region where its distribution is independent of $\mathbf{s}_1$. If not, it discards everything, draws a fresh $\mathbf{y}$, and retries. Only a leak-free $\mathbf{z}$ is ever released.

\begin{figure}[H]
\centering
\begin{tikzpicture}[node distance=5.5mm and 11mm]
  \node[flownode] (y) {sample mask $\mathbf{y}$};
  \node[flownode, right=of y] (w) {$\mathbf{w}=A\mathbf{y}$};
  \node[flownode, right=of w] (c) {$\tilde c=H(\mu\,\|\,\mathbf{w}_1)$};
  \node[flownode, below=of c] (z) {$\mathbf{z}=\mathbf{y}+c\,\mathbf{s}_1$};
  \node[diamond, draw=pqcorange, thick, fill=white, align=center,
        inner sep=1pt, font=\scriptsize, left=of z] (chk) {$\mathbf{z}$\\safe?};
  \node[keynode, left=of chk] (out) {output\\[-2pt]\scriptsize $\sigma=(c,\mathbf{z},\mathbf{h})$};
  \draw[flowarrow] (y) -- (w);
  \draw[flowarrow] (w) -- (c);
  \draw[flowarrow] (c) -- (z);
  \draw[flowarrow] (z) -- (chk);
  \draw[flowarrow] (chk) -- node[above,font=\scriptsize,pqcteal]{yes} (out);
  \draw[flowarrow] (chk.north) |- ([yshift=3.5mm]y.north)
        node[near end,above,font=\scriptsize,pqcorange]{no: restart} -- (y.north);
\end{tikzpicture}
\caption{The ML-DSA signing loop. A fresh mask $\mathbf{y}$ is drawn each round; the response $\mathbf{z}$ is released only when the safety checks confirm its distribution is independent of $\mathbf{s}_1$. On average only a few iterations are needed. The hint $\mathbf{h}$ lets the verifier recover the high bits despite the discarded low part $\mathbf{t}_0$.}
\label{fig:mldsa-abort}
\end{figure}

\section{Signing and verification}

Bringing the phases together gives the signing algorithm. Two bounds mark the safe region: $\|\mathbf{z}\|_\infty$ must stay below $\gamma_1-\beta$, and the low bits of $\mathbf{w}-c\mathbf{s}_2$ must stay below $\gamma_2-\beta$. A hint $\mathbf{h}$ compensates for the omitted $\mathbf{t}_0$.

\begin{algorithm}[H]
\caption{ML-DSA.Sign (teaching form)}
\label{alg:mldsa-sign}
\begin{algorithmic}[1]
\Require Secret key $(\rho,K,tr,\mathbf{t}_0,\mathbf{s}_1,\mathbf{s}_2)$, message $M$, 32-byte value $rnd$
\Ensure Signature $\sigma=(c,\mathbf{z},\mathbf{h})$
\State $A \gets \textsc{ExpandA}(\rho)$;\quad $\mu \gets H(tr\,\|\,M')$
\State $\kappa \gets 0$;\quad $\rho'' \gets H(K\,\|\,rnd\,\|\,\mu)$ \Comment{$rnd=0$ gives the deterministic option}
\Loop
  \State $\mathbf{y} \gets \textsc{ExpandMask}(\rho'', \kappa)$;\quad $\kappa \gets \kappa+\ell$ \Comment{fresh mask each round}
  \State $\mathbf{w} \gets A\mathbf{y}$;\quad $\mathbf{w}_1 \gets \textsc{HighBits}(\mathbf{w})$
  \State $\tilde c \gets H(\mu\,\|\,\mathbf{w}_1)$;\quad $c\gets\textsc{SampleInBall}(\tilde c)$
  \State $\mathbf{z} \gets \mathbf{y} + c\,\mathbf{s}_1$
  \If{$\|\mathbf{z}\|_\infty \ge \gamma_1-\beta$ \textbf{ or } $\|\textsc{LowBits}(\mathbf{w}-c\mathbf{s}_2)\|_\infty \ge \gamma_2-\beta$}
    \State \textbf{continue} \Comment{unsafe: discard and resample}
  \EndIf
  \State $\mathbf{h} \gets \textsc{MakeHint}(c\mathbf{t}_0,\ \mathbf{w}-c\mathbf{s}_2+c\mathbf{t}_0)$
  \State \Return $(c,\mathbf{z},\mathbf{h})$
\EndLoop
\end{algorithmic}
\end{algorithm}

Verification recomputes the commitment from public data and checks the challenge hash matches, while also confirming $\mathbf{z}$ is genuinely short -- the short-vector check is what confronts a forger with Module-SIS.

\begin{algorithm}[H]
\caption{ML-DSA.Verify (teaching form)}
\label{alg:mldsa-verify}
\begin{algorithmic}[1]
\Require Public key $(\rho,\mathbf{t}_1)$, message $M$, signature $(c,\mathbf{z},\mathbf{h})$
\Ensure \emph{accept} or \emph{reject}
\State $A \gets \textsc{ExpandA}(\rho)$;\quad $\mu \gets H(H(\rho\,\|\,\mathbf{t}_1)\,\|\,M)$
\State $\mathbf{w}_1' \gets \textsc{UseHint}\big(\mathbf{h},\ A\mathbf{z} - c\,\mathbf{t}_1\cdot 2^d\big)$
\State $c' \gets H(\mu\,\|\,\mathbf{w}_1')$
\If{$c'=c$ \textbf{ and } $\|\mathbf{z}\|_\infty < \gamma_1-\beta$}
  \State \Return \emph{accept}
\Else
  \State \Return \emph{reject}
\EndIf
\end{algorithmic}
\end{algorithm}

Because $A\mathbf{z} - c\,\mathbf{t} = A\mathbf{y} - c\mathbf{s}_2$ agrees with $\mathbf{w}=A\mathbf{y}$ in its high bits, and the hint repairs the gap left by publishing only $\mathbf{t}_1$, both sides hash to the same challenge exactly when the signature is genuine.

\section{Why a forger is stuck}

An attacker without $sk$ may freely choose $\mathbf{z}$ and $c$, but to pass verification they must satisfy $c = H(\mu\,\|\,\textsc{HighBits}(A\mathbf{z}-c\mathbf{t}))$ \emph{and} keep $\mathbf{z}$ short. Because $H$ behaves like a random function, its output cannot be steered; the only route is to find a genuinely short $\mathbf{z}$ consistent with the equation -- a Module-SIS instance. And they cannot instead learn $\mathbf{s}_1$ from public data, because that is Module-LWE. The two hard problems shut both doors at once.

\section{SageMath experiment}

The core cancellation is easy to verify in miniature. Using small parameters, sample a key, form a masked response, and confirm the verification identity holds before any signature is even released.

\begin{lstlisting}[language=Python]
# Toy sizes; real ML-DSA uses n=256, q=8380417.
from random import Random
from sage.all import GF, Matrix, PolynomialRing, vector

q = 8380417
R.<x> = PolynomialRing(GF(q))
Rq = R.quotient(x^8 + 1)  # toy ring
rng = Random("mldsa-cancellation")

def small_poly():
    return Rq([rng.choice([-1, 0, 1]) for _ in range(8)])

def uniform_poly():
    return Rq([rng.randrange(q) for _ in range(8)])

A = Matrix(Rq, 2, 2, [uniform_poly() for _ in range(4)])
s1 = vector(Rq, [small_poly(), small_poly()])
s2 = vector(Rq, [small_poly(), small_poly()])
t = A * s1 + s2
y = vector(Rq, [uniform_poly(), uniform_poly()])
c = small_poly()
# For a mask y and challenge c, check the verifier's identity:
#   A*z - c*t  ==  A*y - c*s2   (should hold exactly)
z = y + c*s1
assert A*z - c*t == A*y - c*s2
\end{lstlisting}

Then repeat with several masks and count how often $\|\mathbf{z}\|_\infty$ exceeds the bound: this is exactly the rejection rate, and it shows why a signature takes a small, variable number of attempts.

\section{A complete SageMath implementation}
\label{sec:fips204-impl}

The cancellation check above is a fragment. The companion file
\sagefile{fips204\_mldsa.sage} is the whole standard: all forty-nine
numbered algorithms of FIPS~204, each as its own function annotated with its
algorithm number, for ML-DSA-44, 65, and 87. That includes the parts this
chapter has treated as black boxes -- \textsc{ExpandA}, \textsc{ExpandS},
\textsc{ExpandMask}, \textsc{SampleInBall}, \textsc{Power2Round},
\textsc{Decompose}, \textsc{MakeHint}, \textsc{UseHint} -- as well as the bit
packing, the hint encoding, the pre-hash variants \textsc{HashML-DSA}, and
Montgomery reduction.

As in FIPS~203 the algebra is carried by Sage. Here, though, the standard is
careful to distinguish the ring $R_q$ from the ring $R$ of integer
polynomials with a restricted coefficient range, and so is the code:
quantities reduced modulo $q$ (the secrets $\mathbf{s}_1,\mathbf{s}_2$, the
mask $\mathbf{y}$, the response $\mathbf{z}$) live in the Sage quotient ring
$R_q = \mathbb{F}_{8380417}[X]/(X^{256}+1)$, while $\mathbf{t}_1$,
$\mathbf{w}_1$, $\mathbf{r}_0$, and the hint $\mathbf{h}$ stay integer lists.
The NTT domain $T_q$ is the direct product $(\mathbb{Z}_q)^{256}$, so
\textsc{MultiplyNTT} (Algorithm~45) is literally a coordinatewise product.

The rejection loop of Algorithm~\ref{alg:mldsa-sign} then reads almost exactly
as the standard writes it:

\begin{lstlisting}[language=Python]
kappa = 0
while True:
    y  = self.ExpandMask(rhopp, kappa)
    yh = [NTT(Rq(p)) for p in y]
    w  = [NTTInverse(a) for a in MatrixVectorNTT(Ah, yh)]
    w1 = HighBitsVec(w, self.gamma2)

    ctilde = H(mu + self.w1Encode(w1), self.lam // 4)
    c      = self.SampleInBall(ctilde)
    ch     = NTT(Rq(c))

    cs1  = [NTTInverse(MultiplyNTT(ch, a)) for a in s1h]
    cs2  = [NTTInverse(MultiplyNTT(ch, a)) for a in s2h]
    z    = [Rq(y[i]) + cs1[i] for i in range(self.ell)]
    wcs2 = [w[i] - cs2[i] for i in range(self.k)]
    r0   = LowBitsVec(wcs2, self.gamma2)

    kappa += self.ell
    if inf_norm_vec(z) >= self.gamma1 - self.beta:      # unsafe: retry
        continue
    if inf_norm_ints(r0) >= self.gamma2 - self.beta:
        continue

    ct0 = [NTTInverse(MultiplyNTT(ch, a)) for a in t0h]
    h   = MakeHintVec([-p for p in ct0],
                      [wcs2[i] + ct0[i] for i in range(self.k)],
                      self.gamma2)
    if inf_norm_vec(ct0) >= self.gamma2:
        continue
    if sum(sum(p) for p in h) > self.omega:
        continue

    return self.sigEncode(ctilde, [centered(zi) for zi in z], h)
\end{lstlisting}

Every \texttt{continue} is one of the discarded attempts of
Figure~\ref{fig:mldsa-abort}; counting them over many signatures reproduces
the ``repetitions'' row of the parameter table, roughly four for ML-DSA-44.
Note also that $\kappa$ advances before the checks, so each retry draws a
genuinely fresh mask.

\paragraph{Checking the transform against the algebra.}
The function \texttt{verify\_ntt\_is\_a\_ring\_isomorphism()} at the end of
\sagefile{fips204\_mldsa.sage} -- a namesake of the ML-KEM check of
Chapter~\ref{ch:fips203}, but over a different ring -- uses Sage to
confirm that $\mathrm{NTT}^{-1}$ inverts $\mathrm{NTT}$; that the $j$-th NTT
coefficient is the evaluation $w(\zeta^{2\,\mathrm{BitRev8}(j)+1})$, which is
what makes $T_q$ a direct product; and that \textsc{MultiplyNTT} transports
multiplication in $R_q$ as Sage computes it. It also checks that
\textsc{MontgomeryReduce} really returns $a\cdot 2^{-32} \bmod q$.

\paragraph{Validation.}
The implementation is checked against NIST's ACVP vectors~\cite{acvp} for FIPS~204 --
key generation, signature generation, and signature verification -- and
reproduces them byte for byte. The signature-generation vectors are the
demanding ones, because they cover every combination the standard permits:
deterministic and hedged signing, the pure and pre-hash message encodings
across all twelve approved pre-hash functions, the external and internal
signing interfaces, and the external-$\mu$ mode in which the message
representative is computed in a separate cryptographic module. Run
\texttt{make kat} from the \texttt{latex/} directory.

\section{Pitfalls}

\paragraph{Pitfall 1: thinking the mask can be reused or made small.}
The mask $\mathbf{y}$ must be fresh every attempt and large enough to hide $c\,\mathbf{s}_1$. Reusing a mask, or shrinking it to avoid rejections, leaks the secret. This is the whole reason for the abort loop.

\paragraph{Pitfall 2: thinking the abort loop is a timing leak about the secret.}
The number of iterations depends on the mask, which is fresh randomness, not on the secret in an exploitable way. FIPS~204 derives the mask seed from $K\,\|\,rnd\,\|\,\mu$: choosing a zero $rnd$ gives deterministic signing, while a fresh $rnd$ supports randomized or hedged signing.

\paragraph{Pitfall 3: thinking $A$ is stored in the key.}
As in ML-KEM, $A$ is regenerated from $\rho$ every time. The public key is only $(\rho,\mathbf{t}_1)$.

\section{Parameter sets}

\begin{table}[H]
\centering
\begin{tabular}{lccc}
\toprule
Scheme & $(k,\ell)$ & NIST level & Comparable to \\
\midrule
ML-DSA-44 & $(4,4)$ & 2 & AES-128 \\
ML-DSA-65 & $(6,5)$ & 3 & AES-192 \\
ML-DSA-87 & $(8,7)$ & 5 & AES-256 \\
\bottomrule
\end{tabular}
\caption{The three standardized ML-DSA parameter sets. All share $n=256$ and $q=8380417$.}
\end{table}

\section{Summary}

ML-DSA reuses the module-lattice world of Kyber and adds one genuinely new ingredient, Fiat--Shamir with aborts:

\begin{itemize}
  \item key generation is a deterministic expansion from one seed $\xi$, exactly like FIPS 203;
  \item a signature is a commitment--challenge--response transcript made non-interactive by hashing;
  \item the response $\mathbf{z}=\mathbf{y}+c\,\mathbf{s}_1$ carries the secret, so it is released only when rejection sampling confirms it hides the secret completely;
  \item unforgeability rests on Module-LWE (hiding the key) and Module-SIS (blocking short-vector forgeries).
\end{itemize}

The next chapter stays with lattices but takes a very different route to a signature: FN-DSA (Falcon), which abandons Fiat--Shamir for a hash-and-sign construction over NTRU lattices and, in exchange, produces the most compact signatures of the three standards.
